\section{Einops 与 Shape：让维度变成可读账本}

PyTorch 的 `transpose(-2, -1)`、`sum(dim=-1)` 很强，但隐藏了维度语义。Einops 用命名维度写 tensor operation，让表达式更接近数学。

\begin{lstlisting}[caption={传统 PyTorch 写法容易隐藏维度含义}]
x = torch.ones(2, 2, 3)      # batch seq hidden
y = torch.ones(2, 2, 3)      # batch seq hidden
z = x @ y.transpose(-2, -1)  # batch seq seq
\end{lstlisting}

\begin{lstlisting}[caption={Einops/einsum 用名字表达维度收缩}]
z = einsum(
    x, y,
    "batch seq1 hidden, batch seq2 hidden -> batch seq1 seq2"
)
\end{lstlisting}

\begin{knowledgebox}{术语消化：einsum、reduce、rearrange}
\begin{itemize}
    \item \textbf{einsum}：按名字指定哪些维度相乘、哪些维度求和、哪些维度保留。
    \item \textbf{reduce}：沿某些维度做 sum/mean/max/min。
    \item \textbf{rearrange}：拆分、合并、重排维度，例如把 \texttt{total\_hidden} 拆成 \texttt{heads hidden}。
\end{itemize}
\end{knowledgebox}

\begin{warningbox}{Shape 错误是隐性 bug}
维度顺序错了，代码可能仍能运行，只是语义错。命名维度能减少这类错误，也为 tensor parallelism、sequence parallelism 和 attention head 切分打基础。
\end{warningbox}

\subsection{本章小结}

Einops 代表一种资源核算习惯：每个维度都要有名字。后续计算 FLOPs、activation memory 和 parallelism 时，维度语义比张量本身更重要。

